\documentclass[12pt,a4paper]{article}

\usepackage{fontawesome5}
\providecommand{\faPencilAlt}{\faIcon{pencil-alt}}

\usepackage{styles/couleurs}
\usepackage{styles/mise_en_page}
\usepackage{styles/boites}
\usepackage{styles/maths}
\usepackage{styles/figures}

\usepackage{enumitem}
\usepackage{array}
\usepackage{tabularx}
\usepackage{amsmath}
\usepackage{amssymb}
\usepackage{graphicx}
\usepackage{tikz}
\usepackage{xparse}

\usetikzlibrary{calc,angles,quotes,arrows.meta}

\setlist[itemize]{leftmargin=1.4em,itemsep=0.15em,topsep=0.2em}
\setlist[enumerate]{leftmargin=*,itemsep=0.22em,topsep=0.25em}

\newcommand{\lignepointillee}{\noindent\dotfill\par\vspace{0.45em}}

\newcommand{\lignesreponse}[1]{%
\begin{tcolorbox}[
enhanced,colback=white,colframe=NeutreC!40,boxrule=0.5pt,
arc=1.5mm,left=3mm,right=3mm,top=2mm,bottom=1mm,
before skip=2pt,after skip=2pt]
\foreach \i in {1,...,#1}{\lignepointillee}
\end{tcolorbox}}

\newcommand{\titreseq}[1]{%
\begin{center}
{\sffamily\large\bfseries\color{SavoirC}#1}
\end{center}}

\begin{document}

\titreseq{Activité 2 -- Un même angle, un même multiplicateur ?}

\remarque{Les mesures sont approximatives. Les petits écarts entre groupes sont normaux :
ils viennent de la précision du tracé et de la mesure.}

\section*{Partie A -- Trois familles de triangles}

Chaque groupe étudie une famille de triangles rectangles ayant le même angle aigu.

\begin{center}
\begin{minipage}[t]{0.31\textwidth}
\centering
\textbf{Famille 1 : \(35^\circ\)}\par\medskip
\begin{tikzpicture}[scale=0.62]
  \coordinate (O) at (0,0);
  \draw[line width=0.9pt] (O)--(7.7,0);
  \draw[line width=0.9pt] (O)--({8*cos(35)},{8*sin(35)});
  \foreach \r/\name in {4/A,6/B,8/C}{
    \coordinate (\name) at ({\r*cos(35)},{\r*sin(35)});
    \coordinate (\name-p) at ({\r*cos(35)},0);
    \draw (\name)--(\name-p);
    \draw ($(\name-p)+(-0.28,0)$)--($(\name-p)+(-0.28,0.28)$)--($(\name-p)+(0,0.28)$);
  }
  \draw pic[draw=ActionC,angle radius=8mm,"$35^\circ$"] {angle=C-p--O--C};
\end{tikzpicture}
\end{minipage}
\hfill
\begin{minipage}[t]{0.31\textwidth}
\centering
\textbf{Famille 2 : \(50^\circ\)}\par\medskip
\begin{tikzpicture}[scale=0.62]
  \coordinate (O) at (0,0);
  \draw[line width=0.9pt] (O)--(5.3,0);
  \draw[line width=0.9pt] (O)--({8*cos(50)},{8*sin(50)});
  \foreach \r/\name in {4/D,6/E,8/F}{
    \coordinate (\name) at ({\r*cos(50)},{\r*sin(50)});
    \coordinate (\name-p) at ({\r*cos(50)},0);
    \draw (\name)--(\name-p);
    \draw ($(\name-p)+(-0.28,0)$)--($(\name-p)+(-0.28,0.28)$)--($(\name-p)+(0,0.28)$);
  }
  \draw pic[draw=ActionC,angle radius=8mm,"$50^\circ$"] {angle=F-p--O--F};
\end{tikzpicture}
\end{minipage}
\hfill
\begin{minipage}[t]{0.31\textwidth}
\centering
\textbf{Famille 3 : \(70^\circ\)}\par\medskip
\begin{tikzpicture}[scale=0.62]
  \coordinate (O) at (0,0);
  \draw[line width=0.9pt] (O)--(3,0);
  \draw[line width=0.9pt] (O)--({8*cos(70)},{8*sin(70)});
  \foreach \r/\name in {4/G,6/H,8/I}{
    \coordinate (\name) at ({\r*cos(70)},{\r*sin(70)});
    \coordinate (\name-p) at ({\r*cos(70)},0);
    \draw (\name)--(\name-p);
    \draw ($(\name-p)+(-0.28,0)$)--($(\name-p)+(-0.28,0.28)$)--($(\name-p)+(0,0.28)$);
  }
  \draw pic[draw=ActionC,angle radius=8mm,"$70^\circ$"] {angle=I-p--O--I};
\end{tikzpicture}
\end{minipage}
\end{center}

\section*{Partie B -- Mesurer et calculer}

Pour les trois triangles de ta famille :

\begin{enumerate}
  \item mesure l'hypoténuse ;
  \item mesure le côté de l'angle droit qui touche l'angle étudié ;
  \item Calcule :
  \[
  \text{côté adjacent}\div\text{hypoténuse},
  \]
  puis arrondis au centième.
\end{enumerate}

Famille n° ...

\renewcommand{\arraystretch}{1.55}
\begin{tabularx}{\textwidth}{
|>{\centering\arraybackslash}p{2.3cm}
|>{\centering\arraybackslash}X
|>{\centering\arraybackslash}X
|>{\centering\arraybackslash}X|}
\hline
\textbf{Triangle} & \textbf{Hypoténuse} & \textbf{Côté adjacent} &
\textbf{Adjacent \(\div\) hypoténuse}\\
\hline
1 & & & \\
\hline
2 & & & \\
\hline
3 & & & \\
\hline
4 & & & \\
\hline
5 & & & \\
\hline
6 & & & \\
\hline
\end{tabularx}

\section*{Partie C -- Conjecturer}

\begin{enumerate}
  \item Que remarques-tu dans une même famille ?
  \item Le nombre obtenu dépend-il de la taille du triangle ?
  \item Compare ce nombre avec celui des autres familles.
  \item De quoi ce nombre semble-t-il dépendre ?
\end{enumerate}

\lignesreponse{8}

\section*{Partie D -- Utiliser la conjecture}

Dans une famille de triangles rectangles ayant un angle de \(35^\circ\),
le nombre obtenu est environ \(0{,}82\).

\begin{enumerate}
  \item Si l'hypoténuse mesure \(10\) cm, estime le côté adjacent.
  \item Si l'hypoténuse mesure \(25\) cm, estime le côté adjacent.
  \item Retrouve la logique du calcul effectué.
\end{enumerate}

\lignesreponse{7}

\aretenir{Pour un angle aigu fixé, le quotient
\[
\text{côté adjacent}\div\text{hypoténuse}
\]
semble constant. Ce nombre dépend de l'angle, pas de la taille du triangle.}

\end{document}
